\chapter{Code source étendu}
\label{ann:code_extended}

\section*{F.1 — Endpoint FastAPI d'analyse d'image}

\begin{lstlisting}[language=python,
  caption={Endpoint FastAPI POST /api/v1/agent/analyze-image (extrait simplifié)},
  label={lst:analyze_image}]
@router.post("/analyze-image", response_model=AgentAnalysisResponse)
async def analyze_image(
    file: UploadFile = File(...),
    question: str = Form(default="Qu'est-ce que tu vois?"),
    language: str = Form(default="darija"),
    db: AsyncSession = Depends(get_db),
    current_user: User = Depends(get_current_user),
):
    # 1. Validation et sauvegarde temporaire
    if file.content_type not in ALLOWED_IMAGE_TYPES:
        raise HTTPException(status_code=400, detail="Format non supporté")
    tmp_path = await save_upload_temp(file)

    try:
        # 2. Inférence YOLO
        yolo_result = await run_yolo_inference(tmp_path)
        detections = [
            Detection(label=d["label"], confidence=d["conf"],
                      bbox=d["xyxy"])
            for d in yolo_result["detections"]
        ]

        # 3. Requête LLM avec fallback 3-niveaux
        context = build_rag_context(question, detections)
        answer = await llm_service.generate_with_fallback(
            prompt=question,
            context=context,
            language=language,  # "darija" | "fr" | "ar"
            detections=detections,
        )

        # 4. Journalisation MLflow
        mlflow.log_metric("yolo_detections", len(detections))
        mlflow.log_param("llm_fallback_level", answer.fallback_level)

        return AgentAnalysisResponse(
            detections=detections,
            answer=answer.text,
            confidence=answer.confidence,
            fallback_level=answer.fallback_level,
            processing_time_ms=answer.latency_ms,
        )
    finally:
        tmp_path.unlink(missing_ok=True)
\end{lstlisting}

\section*{F.2 — Hook React \texttt{useNetworkSync}}

\begin{lstlisting}[language=Java,
  caption={Hook React useNetworkSync --- synchronisation offline/online (IndexedDB + API)},
  label={lst:use_network_sync}]
import { useState, useEffect, useCallback } from 'react';
import { db } from '../db/dexie';
import { syncPendingActions } from '../services/api';
import { useToast } from './useToast';

export function useNetworkSync() {
  const [isOnline, setIsOnline] = useState(navigator.onLine);
  const [syncStatus, setSyncStatus] = useState('idle'); // idle | syncing | error
  const { showToast } = useToast();

  const performSync = useCallback(async () => {
    const pending = await db.pendingActions.toArray();
    if (pending.length === 0) return;

    setSyncStatus('syncing');
    try {
      const results = await syncPendingActions(pending);
      const successIds = results
        .filter(r => r.success)
        .map(r => r.localId);
      await db.pendingActions.bulkDelete(successIds);

      showToast(`${successIds.length} action(s) synchronisée(s)`, 'success');
      setSyncStatus('idle');
    } catch (err) {
      setSyncStatus('error');
      showToast('Échec de synchronisation — réessai dans 30s', 'warning');
    }
  }, [showToast]);

  useEffect(() => {
    const handleOnline = () => {
      setIsOnline(true);
      performSync(); // Auto-sync on reconnect
    };
    const handleOffline = () => setIsOnline(false);

    window.addEventListener('online', handleOnline);
    window.addEventListener('offline', handleOffline);
    return () => {
      window.removeEventListener('online', handleOnline);
      window.removeEventListener('offline', handleOffline);
    };
  }, [performSync]);

  return { isOnline, syncStatus, performSync };
}
\end{lstlisting}

\section*{F.3 — Service MQTT IoT (backend Python)}

\begin{lstlisting}[language=python,
  caption={Service MQTT — abonnement aux topics IoT et dispatch WebSocket},
  label={lst:mqtt_service}]
import json
import logging
from aiomqtt import Client as MqttClient
from app.core.config import settings
from app.services.alert_service import evaluate_and_dispatch_alert

logger = logging.getLogger(__name__)

MQTT_TOPICS = {
    "smartfarm/node_a/#": "irrigation",   # Node A: soil, pressure, flow, temp, alert
    "smartfarm/node_b/#": "beehive",      # Node B: weight, hive_temp, ext_temp, ext_hum, alert
}
# Topics réels publiés par le firmware ESP32 :
# smartfarm/node_a/{soil,pressure,flow,pump,valve,alert}
# smartfarm/node_b/{weight,hive_temp,ext_temp,ext_hum,alert}

async def mqtt_listener(websocket_manager):
    async with MqttClient(
        hostname=settings.MQTT_BROKER_HOST,
        port=settings.MQTT_BROKER_PORT,
    ) as client:
        for topic in MQTT_TOPICS:
            await client.subscribe(topic)
            logger.info(f"Abonne au topic MQTT : {topic}")

        async for message in client.messages:
            try:
                payload = json.loads(message.payload.decode())
                topic_str = str(message.topic)

                # Extraction du farm_id depuis le topic (pattern: farm/{id}/...)
                parts = topic_str.split("/")
                farm_id = int(parts[1]) if len(parts) > 1 else None

                # Evaluation des seuils et creation d'alerte si necessaire
                alert = await evaluate_and_dispatch_alert(
                    farm_id=farm_id,
                    sensor_type=MQTT_TOPICS.get(topic_str, "unknown"),
                    payload=payload,
                )

                # Broadcast WebSocket vers tous les clients du tenant
                if alert:
                    await websocket_manager.broadcast_to_farm(
                        farm_id=farm_id,
                        event={"type": "ALERT", "data": alert.dict()},
                    )

            except (json.JSONDecodeError, KeyError, ValueError) as exc:
                logger.warning(f"Payload MQTT invalide sur {message.topic}: {exc}")
\end{lstlisting}

% ============================================================
%  ANNEXE G : Données brutes et formats
% ============================================================
